% Источники: exam/materials/моделирование.pdf, разделы 26--27; https://github.com/HumsterProgrammer/iu9-notes/blob/main/8sem/Моделирование/Моделирование_лекция-04_19-02-26.md.
\section{Подход к имитации на основе активностей и событий.}

\subsection*{Подход на основе активностей}

\textbf{Активность} — элементарное действие, характеризующее поведение одного элемента системы за конечный промежуток модельного времени.

Активность описывается кортежем:
\[
  A_{ij} = \bigl\{Al_{ij},\; \Delta\tau_{ij}\bigr\},
\]
где $Al_{ij}$ -- алгоритм обработки, а $\Delta\tau_{ij}$ -- квант модельного времени, на который меняется модельное время объекта или компоненты.

Работа управляющей программы:
\begin{enumerate}
  \item Фиксируется текущий момент модельного времени $\tau$.
  \item Для каждого элемента перебираются все активности; активности, момент начала которых совпадает с $\tau$, выполняются.
  \item Модельное время продвигается на шаг $\Delta\tau$.
\end{enumerate}

Подход применяется для систем с небольшим количеством компонент и состояний, обеспечивает высокую точность моделирования и прост в организации.

\subsection*{Подход на основе событий}

\textbf{Событие} — мгновенное изменение состояния системы в момент времени $\tau$. В лекционных конспектах событие описывается как конкретное значение времени наступления события:
\[
  s_{ij} \in S_i = \bigl\{s^*_{ij},\; \tau_{ij}\bigr\},
\]
где $s_{ij}$ -- $j$-е событие $i$-й компоненты, $s^*_{ij}$ -- его упрощенная формализация, а $\tau_{ij}$ -- показатель модельного времени.

Работа управляющей программы:
\begin{enumerate}
  \item Все запланированные события хранятся в \textbf{календаре событий} (очередь с приоритетом по времени).
  \item Извлекается событие с наименьшим временем $\tau_{\min}$.
  \item Модельное время устанавливается $\tau = \tau_{\min}$.
  \item Выполняется обработчик события; при необходимости в календарь добавляются новые события.
\end{enumerate}

Подход применяется при составлении расписаний, если для решения задачи не требуется более подробное описание состояний системы. Пример: расписание занятий, где важны значения признаков <<где>>, <<когда>>, <<кто>> и <<занятие состоялось/не состоялось>>.

\clearpage
